\documentclass[10pt,a4paper]{amsart}
\usepackage[margin=19mm]{geometry}
\usepackage{amsmath,amssymb,mathtools}
\usepackage[scr=rsfs,cal=euler]{mathalfa}
\usepackage{xfrac}
\usepackage[unicode,hidelinks,bookmarks=false]{hyperref}
\newcommand{\ie}{i.e.\ }
\newcommand{\cf}{cf.\ }
\newcommand{\resp}{respectively\ }
\newcommand{\Subsection}{\subsection}
\providecommand{\nobreakdash}{\nobreak\hbox{-}}
\setlength{\emergencystretch}{2em}
\multlinegap0pt
\usepackage{bm}
\usepackage{bbm}
\usepackage{dsfont}
\usepackage{xspace}
\usepackage{cancel}
\usepackage{mathtools}
\usepackage{amsthm}
\makeatletter
\@ifundefined{lemma}{\newtheorem{lemma}{Lemma}}{}
\makeatother
\title[Circles and line segments as independence attractors of grap]{Circles and line segments as independence attractors of graphs}
\author{Garima Khetawat, Moumita Manna, Tarakanta Nayak}
\date{Source: arXiv:2505.20898v1 (CC BY 4.0)}
\begin{document}
\maketitle

\noindent\emph{Source and licence.} Circles and line segments as independence attractors of graphs. Version arXiv:2505.20898v1 (CC BY 4.0), distributed under the Creative Commons Attribution 4.0 International licence, as recorded in the arXiv licence metadata for this exact version and shown on the abstract page https://arxiv.org/abs/2505.20898v1. Full rights notice: licenses/arxiv-2505.20898-CC-BY-4.0.txt.\par\smallskip

\noindent\textit{Editorial scope.} Counted unit: the Lemma whose source label is \texttt{Nograph-k5-key} in the article (source0.tex lines 434--437, source label \texttt{Nograph-k5-key}), delivered together with the article's own complete original proof of it (source0.tex lines 438--459, one \texttt{proof} environment of 22 lines). No mathematics is written, repaired or re-derived: statement and proof are the article's own text line for line. Required context retained uncounted: none. Every definition, display and label of those lines is retained unchanged. Omitted from this package and not redistributed: 1--433, 460--1766. Nothing inside a retained excerpt is omitted or altered: every retained body is delivered exactly as the article's lines are written, so no omission receipt of any kind accompanies this package. Citations: the retained excerpts carry no citation command at all, so no bibliography entry of the article is transcribed and no reference is redistributed. Reference adaptation: none was needed and none was made; every reference inside the delivered unit resolves inside it, so no unresolved reference or citation is produced. Printed slips kept literal: no printed word, symbol, hypothesis or number of the counted statement or of its proof was altered. Layout and class: the article's own class is \texttt{article} with packages including amsfonts, amssymb, amscd, amsmath, enumerate, verbatim, calc, graphicx, hyperref, babel, geometry, leftindex, setspace, tikz; the wrapper is the lane's pinned \texttt{amsart} editorial wrapper at 10pt on A4, which restates the theorem-like environments the retained text needs and adds the article's own macro definitions that the retained text needs (none), with extra wrapper packages (bm, bbm, dsfont, xspace, cancel, mathtools). The delivered numbering is the unit's own. Assets: no retained span contains a figure environment, a graphics-inclusion command, a PDF-page inclusion, a table or a TikZ picture; no image file is copied into this package, the package declares no asset dependency and no image file is redistributed with it. Licence: version arXiv:2505.20898v1 of this article is distributed under the Creative Commons Attribution 4.0 International licence, as recorded in the arXiv licence metadata for this exact version and shown on the abstract page https://arxiv.org/abs/2505.20898v1. A full rights notice is delivered at licenses/arxiv-2505.20898-CC-BY-4.0.txt. Characters: every retained excerpt is pure ASCII, so the delivered engine, XeLaTeX with its default Latin Modern text font, reports no missing character.
  \begin{lemma}
  	Let $G$ be a simple graph with at least two vertices. Also, let $N, E$ and $T$ denote the number of vertices, edges and triangles in $G$ respectively. If $4E > N^2$ then $T \geq \frac{E}{3N} (4E -N^2)$.
  	\label{Nograph-k5-key}
  \end{lemma}

  \begin{proof}
For an edge $e$ joining the vertices $v_1, v_2$, let $n_0 (e), n_1 (e)$ and $n_2 (e) $ be the number of vertices different from $v_1, v_2$ that are adjacent to none of, exactly one of and both of $v_1, v_2$, respectively. Then  \begin{equation}
	n_0 (e)+ n_1 (e)+n_2 (e) =N-2  ~\mbox{for every edge}~ e.
	\label{N-2}
\end{equation}
 Let $P_n$ denote the path  graph on $n$ vertices. We call it  an $n$-path. Let $T_1$ be the number of times  $P_2 \cup K_1 $, the disjoint union of $P_2$ and $K_1$,  occurs in $G$ as an induced subgraph. Then $\sum_{e \in E(G)} n_0 (e)=T_1$.
 Similarly, let $T_2$ and $T $ denote the number of $3$-paths in $G$ that is not a subgraph of any triangle, and triangles in $G$ respectively.  Taking the sum over all the edges of $G$, we get  $\sum_{e \in E(G)}n_1 (e) =2 T_2$ and  $\sum_{e \in E(G)}n_2 (e) = 3T$.  To see these equalities, note that each $3$-path in $G$ is counted exactly twice and each triangle is counted exactly thrice while the sum over all edges is taken. Thus, $ T_1 +2 T_2+3T  = \sum_{e \in E(G)}n_0 (e) + n_1 (e)+ n_2 (e)$ and using Equation~(\ref{N-2}) we get 
 \begin{equation}
T_1 + 2 T_2 +3 T  =(N-2)E.
\label{T1T2T3} 
\end{equation} 
Consider a vertex $v$ of $G$ with degree $\deg(v)$ at least $2$. For each choice of two edges incident up on $v$,   either there is a $3$-path (that is not a subgraph of any triangle) with $v$ as its  middle vertex, or there is a triangle containing $v$. There are $\binom{\deg(v)}{2}$ such distinct choices.  On the other hand, each $3$-path that is not a subgraph of any triangle is counted exactly once and each triangle is counted exactly thrice when all such choices are considered for all vertices with degree at least $2$. This gives that $	T_2 +3 T =\sum_{v \in V(G), ~\deg(v) \geq 2} \binom{\deg(v)}{2}.$    Since $ \sum_{v \in V(G), ~\deg(v) < 2} \binom{\deg(v)}{2}=0, $  $ \sum_{v \in V(G), ~\deg(v) \geq 2} \binom{\deg(v)}{2}= \sum_{v\in V(G)} \binom{\deg(v)}{2}$. Further, $-2E + \sum_{v\in V(G)}  \deg(v)^2 =-\sum_{v\in V(G)} \deg(v) + \sum_{v\in V(G)}  \deg(v)^2 = 2 \sum_{v\in V(G)}  \binom{\deg(v)}{2}$. In other words, 
 \begin{equation}
 2	T_2 +6 T  =-2 E +\sum_{v \in V(G)} \deg(v)^2.
	\label{T2T3}
\end{equation}
Eliminating $T_2$ from Equations (\ref{T1T2T3})  and (\ref{T2T3}), it is found that \begin{equation}
	3T =T_1 -NE +\sum_{v \in V(G)} \deg(v)^2.
	\label{3T3}
	\end{equation} 
Now, $\sum_{1 \leq i <j\leq N} (x_i -x_j)^2 =N\sum_{1\leq i \leq N} x_i ^2 - (\sum_{1 \leq i \leq N} x_i)^2$ for each set of natural numbers $\{x_1, x_2, \cdots, x_N\}$. Putting $\deg(v_i)=x_i$, observe that $\sum_{1 \leq i \leq N} x_i = 2E$, and consequently,  $N \sum_{v \in V(G)} \deg(v)^2 =4E^2+\sum_{1 \leq i < j\leq N} (\deg(v_i) -\deg(v_j))^2$. Putting this in Equation~(\ref{3T3}) we have $3 N T  = 4E^2 + \sum_{1 \leq i <j \leq N} (\deg(v_i) -\deg(v_j))^2 +NT_1 -N^2 E$. As $ \sum_{1 \leq i <j \leq N} (\deg(v_i) -\deg(v_j))^2 +NT_1 \geq 0$, $3 NT  \geq 4E^2 -N^2 E$, i.e., $$ T   \geq \frac{E(4E -N^2)}{3N}.$$ 
  \end{proof}

\end{document}
